% Current v16 frontier. Stable shape IDs; no existence assertion.
\subsection{残る十二型と，再開する具体的な入口}
\label{ns16:support-agenda}

現在の高分母 support の候補は\emph{十二型・三百六十一標識}である．
二葉付き三角形 $G_{5,3}$ は全排除され，その六十標識を旧十三型から除いた．
以下の $G_{n,j}$ の識別子は維持し，欠番を詰めない．
孤立頂点を省き，$ij$ は無向辺 $\{i,j\}$ を表す．
$\mathbf d$ は降順の次数列，$\beta_1=|E|-n+1$ は cycle rank，
$L$ は葉の数，$N$ は $K_5$ 上の標識付き個数である．

\begin{center}
\footnotesize
\setlength{\tabcolsep}{3pt}
\begin{tabular}{crclrrr}
型 & $|E|$ & $\mathbf d$ & $E$ & $\beta_1$ & $L$ & $N$\\\hline
$G_{4,3}$ & 6 & $(3333)$ & $\{01,02,03,12,13,23\}$ & 3 & 0 & 5 \\
$G_{5,5}$ & 6 & $(33222)$ & $\{01,02,03,12,14,34\}$ & 2 & 0 & 60 \\
$G_{5,7}$ & 6 & $(33321)$ & $\{01,02,03,12,13,24\}$ & 2 & 1 & 60 \\
$G_{5,9}$ & 6 & $(43221)$ & $\{01,02,03,04,12,13\}$ & 2 & 1 & 60 \\
$G_{5,10}$ & 7 & $(33332)$ & $\{01,02,03,12,13,24,34\}$ & 3 & 0 & 30 \\
$G_{5,11}$ & 7 & $(43322)$ & $\{01,02,03,04,12,13,24\}$ & 3 & 0 & 60 \\
$G_{5,12}$ & 7 & $(43331)$ & $\{01,02,03,04,12,13,23\}$ & 3 & 1 & 20 \\
$G_{5,13}$ & 7 & $(44222)$ & $\{01,02,03,04,12,13,14\}$ & 3 & 0 & 10 \\
$G_{5,14}$ & 8 & $(43333)$ & $K_5-\{14,23\}$ & 4 & 0 & 15 \\
$G_{5,15}$ & 8 & $(44332)$ & $K_5-\{24,34\}$ & 4 & 0 & 30 \\
$G_{5,16}$ & 9 & $(44433)$ & $K_5-\{34\}$ & 5 & 0 & 10 \\
$G_{5,17}$ & 10 & $(44444)$ & $K_5$ & 6 & 0 & 1 \\
\end{tabular}
\end{center}

全 $1024$ 部分 graph と自己同型群を
\nolinkurl{verify_support_agenda_v16.py} で再検算している．
表は \nolinkurl{support_agenda_verification_v16.json} の \texttt{shapes} から再生成できる．
非空の support は三角形を含み，$\beta_1\ge2$，辺数六以上である．
四使用頂点の場合は $K_4$ のみである．表は必要条件を通過した候補であり，
整数接着や atlas の同型類の数ではない．

\paragraph{第一候補：$G_{4,3}=K_4$ の二つの残存分岐．}
第五の全低 core を hub とする．共通低 puncture に関する通常辺と例外辺を区別する．
全六辺が例外となる配置は排除され，例外辺数は二または四である．
二例外なら正負の例外分母は $(3,7)$ または $(5,9)$ であり，
通常四辺の graph は $C_4$ または三角形に葉辺を付けた形となる．
固定した node 行列と通常 reverse の下で，通常辺の流れの自由度は高々一つの
cycle circulation である．これに三高分母の残る pair equations と高辺 clutching を課す．
四例外なら通常二辺の分母は二十五以下，その二辺が端点を共有しないなら七以下である．
例外分母の上界はまだ得られていない．両分岐とも整数接着の存在・不存在は未解決である．

\paragraph{次候補：$G_{5,5}$ と $G_{5,10}$ の二価 core．}
$G_{5,5}$ の二価頂点 $3,4$ は，高分母近傍の間が低辺なので正の枝に限られる．
頂点 $2$ の近傍 $0,1$ は高辺で結ばれるため，その負の枝を先に捨てない．
高辺と共通低 puncture の組 $(ij;k)$ は
\[
 (02;4),\qquad(12;3),\qquad(34;2)
\]
であり，通常接着なら reverse の絶対値は四以下である．
最大約数に対応する例外接着を残して，三つの二高分母 core の低 row と node を接続する．
$G_{5,10}$ では二価頂点 $4$ が正の枝に限られ，
共通低 puncture の組は $(01;4)$ である．
一般の node 補題は中心の高分母次数や対向 $D$ の相殺を仮定しないので，
三高分母中心が未分類であることを理由にこの制約を保留する必要はない．

\paragraph{他の型と局所理論．}
三高分母の正符号 graph が常に matching となるか，
全ての高 $D$ が非負となり得るかは引き続き未解決である．
特異な rank 二 Schur 系を除外せず，既知の固定 slice の無限族も保持する．
$G_{5,14}=K_{2,2,1}$ では dense mod 五の二模型の整数 lifting が別の入口である．
高分母 support $H=\{ij:q_{ij}>1\}$ と素数 support $G_\ell=\{ij:\ell\mid q_{ij}\}$ を区別する．
局所整数族の存在から五 core の同時接着を結論しない．

二葉付き三角形の十模型・五十二組と $B$ 表現の七組は解決済みであり，再探索しない．
前者と後者は重複するので個数を加算しない．
そこで得た共通低 puncture の node 上界，単位 entry の $(r,s,E)$ 復元，
有限な node 配分から中心三行を復元する手順を，残る型へ再利用する．
優先順位は固定せず，数学的な面白さ，インパクト，幾何的構成への寄与に応じて変更する．
